\subsection{CONJUGACY OF CARTAN SUBALGEBRAS OF SOLVABLE LIE ALGEBRAS}

Let g be a solvable Lie algebra. Denote by \(\mathcal{C}^{\infty}(\mathfrak{g})\) the intersection of the terms of the descending central series of g (Chap. I, §1, no. 5). This is a characteristic ideal of g, and is the smallest ideal m of g such that g/m is nilpotent. Since \(\mathcal{C}^{\infty}(\mathfrak{g}) \subset [\mathfrak{g}, \mathfrak{g}]\) , \(\mathcal{C}^{\infty}(\mathfrak{g})\) is a nilpotent ideal of g (Chap. I, §5, no. 3, Cor. 5 of Th. 1). By Prop. 1 of no. 1, the set of \(e^{ad x}\) , for \(x \in \mathcal{C}^{\infty}(\mathfrak{g})\) , is a subgroup of \(\operatorname{Aut}(\mathfrak{g})\) contained in the group of special automorphisms (Chap. I, §6, no. 8, Def. 6).

THEOREM 3. Let \(\mathfrak{g}\) be a solvable Lie algebra, and let \(\mathfrak{h},\mathfrak{h}'\) be Cartan subalgebras of \(\mathfrak{g}\) . There exists \(x\in \mathcal{C}^{\infty}(\mathfrak{g})\) such that \(e^{\mathrm{ad}x}\mathfrak{h} = \mathfrak{h}'\) .

We argue by induction on \(\dim \mathfrak{g}\) , the case where \(\mathfrak{g} = 0\) being trivial. Let \(\mathfrak{n}\) be a minimal non-zero commutative ideal of \(\mathfrak{g}\) . Let \(\varphi : \mathfrak{g} \to \mathfrak{g} / \mathfrak{n}\) be the canonical morphism. Then \(\varphi(\mathcal{C}^{\infty} \mathfrak{g}) = \mathcal{C}^{\infty}(\mathfrak{g} / \mathfrak{n})\) (Chap. I, §1, no. 5, Prop. 4). Since \(\varphi(\mathfrak{h})\) and \(\varphi(\mathfrak{h}')\) are Cartan subalgebras of \(\mathfrak{g} / \mathfrak{n}\) (§2, no. 1, Cor. 2 of Prop. 4), there exists, by the induction hypothesis, an \(x \in \mathcal{C}^{\infty}(\mathfrak{g})\) such that \(e^{\mathrm{ad} \varphi(x)}\varphi(\mathfrak{h}) = \varphi(\mathfrak{h}')\) . Replacing \(\mathfrak{h}\) by \(e^{\mathrm{ad} x}\mathfrak{h}\) , we can assume that \(\varphi(\mathfrak{h}) = \varphi(\mathfrak{h}')\) , in other words that

\[
\mathfrak {h} + \mathfrak {n} = \mathfrak {h} ^ {\prime} + \mathfrak {n}.
\]

Then \(\mathfrak{h}\) and \(\mathfrak{h}'\) are Cartan subalgebras of \(\mathfrak{h} + \mathfrak{n}\) . If \(\mathfrak{h} + \mathfrak{n} \neq \mathfrak{g}\) , the assertion to be proved follows from the induction hypothesis. Assume from now on that \(\mathfrak{h} + \mathfrak{n} = \mathfrak{h}' + \mathfrak{n} = \mathfrak{g}\) .

By the minimality of \(\mathfrak{n}\) , \([\mathfrak{g},\mathfrak{n}] = \{0\}\) or \([\mathfrak{g},\mathfrak{n}] = \mathfrak{n}\) . If \([\mathfrak{g},\mathfrak{n}] = \{0\}\) , then \(\mathfrak{n} \subset \mathfrak{h}\) and \(\mathfrak{n} \subset \mathfrak{h}'\) (§2, no. 1, Prop. 5), so \(\mathfrak{h} = \mathfrak{h} + \mathfrak{n} = \mathfrak{h}' + \mathfrak{n} = \mathfrak{h}'\) . It remains to consider the case where \([\mathfrak{g},\mathfrak{n}] = \mathfrak{n}\) , so \(\mathfrak{n} \subset \mathcal{C}^{\infty}(\mathfrak{g})\) . The ideal \(\mathfrak{n}\) is a simple \(\mathfrak{g}\) -module; since \(\mathfrak{g} = \mathfrak{h} + \mathfrak{n}\) , and since \([\mathfrak{n},\mathfrak{n}] = \{0\}\) , it follows that \(\mathfrak{n}\) is a simple \(\mathfrak{h}\) -module. If \(\mathfrak{h} \cap \mathfrak{n} \neq \{0\}\) , then \(\mathfrak{n} \subset \mathfrak{h}\) , so \(\mathfrak{g} = \mathfrak{h}\) and \(\mathfrak{h}' = \mathfrak{h}\) . Assume now that \(\mathfrak{h} \cap \mathfrak{n} = \{0\}\) . Then \(\mathfrak{g} = \mathfrak{h} \oplus \mathfrak{n}\) and hence \(\mathfrak{g} = \mathfrak{h}' \oplus \mathfrak{n}\) , since \(\mathfrak{h}\) and \(\mathfrak{h}'\) have the same dimension.

For all \(x \in h\) , let \(f(x)\) be the unique element of n such that \(x - f(x) \in \mathfrak{h}'\) ; if \(x, y \in h\) ,

\[
[ x, y ] - [ x, f (y) ] - [ f (x), y ] = [ x - f (x), y - f (y) ] \in \mathfrak {h} ^ {\prime},
\]

so \(f([x,y]) = [x,f(y)] + [f(x),y]\) . By §1, no. 3, Cor. of Prop. 9, there exists \(a \in \mathfrak{n}\) such that \(f(x) = [x,a]\) for all \(x \in \mathfrak{h}\) . We have \((\operatorname{ad}a)^2(\mathfrak{g}) \subset (\operatorname{ad}a)(\mathfrak{n}) = 0\) , so, for all \(x \in \mathfrak{h}\) ,

\[
e ^ {\operatorname{ad} a} x = x + [ a, x ] = x - f (x).
\]

Thus \(e^{\mathrm{ad}a}(\mathfrak{h}) = \mathfrak{h}'\) . Since \(a \in \mathcal{C}^{\infty}(\mathfrak{g})\) , this completes the proof.

Lemma 3. Let g be a Lie algebra, r its radical, \(\varphi\) the canonical homomorphism from g to g/r, v an elementary automorphism of g/r. There exists an elementary automorphism u of g such that \(\varphi \circ u = v \circ \varphi\) .

We can assume that \(v\) is of the form \(e^{\mathrm{ad}b}\) , where \(b \in \mathfrak{g} / \mathfrak{r}\) and \(\mathrm{ad}b\) is nilpotent. Let \(\mathfrak{s}\) be a Levi subalgebra of \(\mathfrak{g}\) (Chap. I, §6, no. 8, Def. 7) and let \(a\) be the element of \(\mathfrak{s}\) such that \(\varphi(a) = b\) . Since \(\mathrm{ad}_{\mathfrak{s}}a\) is nilpotent, \(\mathrm{ad}_{\mathfrak{g}}a\) is nilpotent (Chap. I, §6, no. 3, Cor. of Prop. 3), and \(u = e^{\mathrm{ad}_{\mathfrak{g}}a}\) is an elementary automorphism of \(\mathfrak{g}\) such that \(\varphi \circ u = v \circ \varphi\) .

PROPOSITION 5. Let g be a Lie algebra, r its radical, h and h' Cartan subalgebras of g, and \(\varphi\) the canonical homomorphism from g to g/r. The following conditions are equivalent:

\begin{enumerate}
\def\labelenumi{(\roman{enumi})}
\item
  \(\mathfrak{h}\) and \(\mathfrak{h}'\) are conjugate by an elementary automorphism of \(\mathfrak{g}\) ;
\item
  \(\varphi (\mathfrak{h})\) and \(\varphi (\mathfrak{h}^{\prime})\) are conjugate by an elementary automorphism of \(\mathfrak{g} / \mathfrak{r}\) .
\item
  \(\Longrightarrow\) (ii): This is clear.
\item
  \(\Longrightarrow\) (i): We assume that condition (ii) is satisfied and prove (i). By Lemma 3, we are reduced to the case where \(\varphi(\mathfrak{h}) = \varphi(\mathfrak{h}')\) . Put \(\mathfrak{k} = \mathfrak{h} + \mathfrak{r} = \mathfrak{h}' + \mathfrak{r}\) , which is a solvable subalgebra of \(\mathfrak{g}\) . Then \(\mathfrak{h}\) and \(\mathfrak{h}'\) are Cartan subalgebras of \(\mathfrak{k}\) , so there exists \(x \in \mathcal{C}^{\infty}(\mathfrak{k})\) such that \(e^{\mathrm{ad}_{\mathfrak{k}}x}\mathfrak{h} = \mathfrak{h}'\) (Th. 3). Since \(\mathfrak{k}/\mathfrak{r}\) is nilpotent, \(\mathcal{C}^{\infty}(\mathfrak{k}) \subset \mathfrak{r}\) ; on the other hand, \(\mathcal{C}^{\infty}(\mathfrak{k}) \subset [\mathfrak{k},\mathfrak{k}] \subset [\mathfrak{g},\mathfrak{g}]\) , so \(x \in \mathfrak{r} \cap [\mathfrak{g},\mathfrak{g}]\) ; by Chap. I, §5, no. 3, Th. 1, \(\mathrm{ad}_{\mathfrak{g}}x\) is nilpotent, so \(e^{\mathrm{ad}_{\mathfrak{g}}x}\) is an elementary automorphism of \(\mathfrak{g}\) transforming \(\mathfrak{h}\) to \(\mathfrak{h}'\) .
\end{enumerate}

